\chapter{Related Works}
\label{cap:trabalhos}

Data leakage has received growing attention as a major threat to the validity,
reproducibility, and practical applicability of ML models. This chapter reviews
the existing literature on the topic. Section~\ref{sec:tr-dl} covers works
dedicated to defining, characterizing, and detecting leakage, as well as those
proposing guidelines to avoid it. Section~\ref{sec:tr-efeitos} covers empirical
studies that measured the consequences of leakage on performance estimates.
Finally, Section~\ref{sec:tr-sumario} summarizes the reviewed works and positions
the contributions of the present study.

\section{Data Leakage}
\label{sec:tr-dl}

The foundational formulation of the problem is due to \citeonline{kaufman2012},
who defined leakage as the introduction of information into the model development
process that would not be legitimately available at prediction time. The authors
formalized the problem, proposed a methodology for its detection, and discussed
avoidance strategies in the context of data mining competitions and predictive
modeling in general.

\citeonline{kapoor2023} expanded this perspective considerably by proposing a
broader taxonomy of leakage mechanisms and demonstrating their role in the
reproducibility crisis affecting ML-based science. Their analysis of studies from
17 scientific fields revealed that methodological flaws involving data
partitioning, pre-processing, and feature engineering are widespread and
frequently result in overly optimistic performance estimates. Their taxonomy is
particularly relevant to the present work, as it explicitly identifies
pre-processing steps applied before data splitting as a distinct and recurrent
leakage mechanism.

The increasing awareness of these risks has motivated the development of
practical guidelines. \citeonline{bernett2024} proposed a set of guiding
questions to help practitioners identify leakage-prone workflows in biological ML
applications, covering both data acquisition and pipeline construction.
\citeonline{apicella2024} discussed common leakage risks in traditional ML and
transfer learning settings, emphasizing how leakage may arise from decisions that
appear methodologically harmless. A classical textbook treatment of the same
issue is provided by \citeonline{hastie2009}, whose discussion of the wrong and
the right way to perform cross-validation describes precisely the mechanism by
which pre-processing applied outside the validation loop contaminates performance
estimates.

Automated detection approaches have also been explored.
\citeonline{drobnjakovic2022} introduced a static analysis framework capable of
detecting leakage patterns directly in ML pipeline code, while
\citeonline{yang2022} proposed methods for identifying leakage-related issues in
computational notebooks, along with process-level recommendations to reduce their
occurrence. These works are complementary to the present study: whereas they aim
to detect whether leakage is present, our objective is to quantify how much the
estimated performance is distorted when it is.

\section{Effects of Data Leakage}
\label{sec:tr-efeitos}

A second group of works investigated the practical consequences of leakage
through empirical experiments, most of them restricted to specific application
domains or to individual pipeline steps.

The earliest and most influential demonstrations come from bioinformatics.
\citeonline{ambroise2002} showed that performing gene selection outside the
cross-validation loop drives the estimated error rate close to zero even on data
containing no real signal, establishing feature selection as a critical leakage
source. \citeonline{varma2006} demonstrated an analogous effect for model
selection, showing that the cross-validation error of a model whose
hyperparameters were themselves chosen by cross-validation is a significantly
biased estimate of the true error, and that nested cross-validation provides a
nearly unbiased alternative.

Subsequent studies confirmed these findings in other domains.
\citeonline{shim2021} showed that performing feature selection before data
partitioning dramatically inflates predictive performance in neuropsychiatric
biomarker discovery, producing misleading conclusions about model effectiveness.
\citeonline{samala2020} reported analogous risks in breast cancer classification
with deep neural networks, where performance rose substantially when feature
selection improperly included the validation set. In medical imaging,
\citeonline{yagis2021} observed substantial inflation in classification accuracy
when MRI slices from the same patient were distributed across training and
testing sets, and \citeonline{bussola2021} showed, in digital pathology, how
correlations between image tiles from the same slide lead to severe performance
overestimation when train-test independence is violated.

Leakage arising from model development procedures rather than from data handling
has also been documented. \citeonline{cawley2010} showed that improper model
selection and hyperparameter optimization procedures can introduce substantial
evaluation bias even when training and testing data sets remain formally
separated, and proposed an unbiased evaluation methodology that accounts for
overfitting occurring during model selection. \citeonline{farrow2019}
demonstrated, in a replication study on discussion forum data, how seemingly
harmless pre-processing decisions may contaminate evaluation protocols and
compromise the validity of reported results.

Particularly relevant to the present work, \citeonline{becker2023} conducted a
systematic empirical assessment of DL effects on classification models,
investigating four leakage scenarios --- normalization, missing values
imputation, feature selection, and hyperparameter tuning --- across 30 PMLB data
sets and six learning algorithms under a unified experimental framework. The
study revealed a consistent pattern of performance overestimation in the presence
of leakage, with the effect being most pronounced for feature selection and
hyperparameter tuning, and with SVM showing comparatively higher susceptibility
than logistic regression. The magnitude of the distortion varied according to
both the learning algorithm and the pipeline step in which leakage originated,
and the data sets could be grouped according to how strongly they responded to
contamination.

\citeonline{becker2023} is the work most closely related to ours, and the direct
motivation for this study. Its experimental design --- paired leakage-free and
leakage-prone pipelines differing only in the presence of contamination --- is
adopted here as well. The present work extends that investigation to regression
tasks, where the target variable is continuous, the applicable algorithms differ,
and the evaluation metrics are unbounded and scale-dependent, none of which can
be assumed to preserve the patterns observed for classification.

\section{Summary}
\label{sec:tr-sumario}

Table~\ref{tab:trabalhos-relacionados} relates the topics addressed by each work
discussed in Sections~\ref{sec:tr-dl} and~\ref{sec:tr-efeitos} to those addressed
by the present study. As can be observed, most previous works concentrate either
on defining and detecting leakage or on assessing its effects in a single
pipeline step or application domain. Only \citeonline{becker2023} evaluated the
four pre-processing and model development steps jointly, and that study is
restricted to classification. To the best of our knowledge, no prior work has
systematically quantified the effects of leakage introduced at multiple pipeline
stages for regression tasks across a diverse collection of benchmark data sets
and learning algorithms, which is the gap this work addresses.

\begin{table}[htb]
\centering
\caption{Summary of the related work and expected contributions of this study.}
\label{tab:trabalhos-relacionados}
\small
\setlength{\tabcolsep}{4pt}
\renewcommand{\arraystretch}{1.15}
\begin{tabular}{l*{10}{c}}
\toprule
\textbf{Work}
 & \begin{sideways}Define / understand DL\end{sideways}
 & \begin{sideways}Propose strategy to avoid DL\end{sideways}
 & \begin{sideways}Detect DL automatically\end{sideways}
 & \begin{sideways}Effects on normalization\end{sideways}
 & \begin{sideways}Effects on value imputation\end{sideways}
 & \begin{sideways}Effects on feature selection\end{sideways}
 & \begin{sideways}Effects on hyperparameter tuning\end{sideways}
 & \begin{sideways}Effects across algorithms\end{sideways}
 & \begin{sideways}Effects across data sets\end{sideways}
 & \begin{sideways}Regression tasks\end{sideways} \\
\midrule
\citeonline{ambroise2002}      &   & X &   &   &   & X &   &   &   &   \\
\citeonline{varma2006}         &   & X &   &   &   &   & X &   &   &   \\
\citeonline{cawley2010}        & X & X &   &   &   &   & X &   &   &   \\
\citeonline{kaufman2012}       & X & X &   &   &   &   &   &   &   &   \\
\citeonline{farrow2019}        &   &   &   &   &   &   &   &   & X &   \\
\citeonline{samala2020}        &   &   &   &   &   & X &   &   &   &   \\
\citeonline{bussola2021}       &   & X &   &   &   &   &   &   & X &   \\
\citeonline{shim2021}          &   &   &   &   &   & X &   &   &   &   \\
\citeonline{yagis2021}         &   &   &   &   &   &   &   & X &   &   \\
\citeonline{drobnjakovic2022}  &   &   & X &   &   &   &   &   &   &   \\
\citeonline{yang2022}          &   & X & X &   &   &   &   &   &   &   \\
\citeonline{kapoor2023}        & X & X &   & X & X & X & X &   & X &   \\
\citeonline{apicella2024}      & X & X &   &   &   &   &   &   &   &   \\
\citeonline{bernett2024}       &   & X &   &   &   &   &   &   &   &   \\
\citeonline{becker2023}        & X &   &   & X & X & X & X & X & X &   \\
\midrule
\textbf{This work}             &   &   &   & X & X & X & X & X & X & X \\
\bottomrule
\end{tabular}
\\[2pt]
\footnotesize Source: The Author
\end{table}
